\section*{Hoofdstuk \ref{ch:b1:radiometric-methods} --- Radiometrische dateringsmethoden}
% values: calc/ch25.py

\begin{solution}{exo:b1:radiometric-methods:1}
(a)~$D^*/P = 2^n - 1 = 3$ geeft $n = 2$: twee halveringstijden. (b)~Rubidium-87:
$e^{\lambda t} - 1 = e^{\num{1.3972e-11} \times \num{3.10e8}} - 1 = 0.004341$;
uranium-235: $e^{\num{9.8485e-10} \times \num{3.10e8}} - 1 = 0.3570$. De
uraniumverhouding is ongeveer 80 keer groter: een gegeven absolute meetfout is een veel
kleiner deel ervan. (c)~$2^{-n} < 0.01$ wanneer $n > \log_2 100 = 6.64$: na ongeveer
zeven halveringstijden.
\end{solution}

\begin{solution}{exo:b1:radiometric-methods:2}
(a)~Koolstof-14 op de houtskool. (b)~Uranium--lood op de zirkoon, of argon--argon
(kalium--argon) op de sanidien, een kaliumrijk \omterm{def:b1:rock-forming-minerals:mineral}{mineraal}. (c)~Uranium--lood op zirkoon,
dat een lange geschiedenis overleeft; rubidium--strontium op hele gesteenten is mogelijk;
kalium--argon loopt het risico van te jonge leeftijden, doordat argon bij latere verhitting verloren ging. (d)~Radiokoolstof op het carbonaat van de schelp, nabij de grens van de methode
(ongeveer \num{70000} jaar). (e)~Geen rechtstreeks: haar leeftijd is omsloten door gedateerde
vulkanische gesteenten erboven en eronder, of gevonden door haar fossielen te correleren met een gedateerde
sectie (\cref{ch:b1:stratigraphic-scale}).
\end{solution}

\begin{solution}{exo:b1:radiometric-methods:3}
(a)~$t = -8033 \ln 0.25 = 8033 \times 1.3863 = \num{11136}$ jaar BP.
(b)~De gemiddelde levensduur is $5700/\ln 2 = \num{8223}$ jaar: $t = 8223 \times 1.3863 =
\num{11400}$ jaar. (c)~Beide nemen aan dat de verhouding koolstof-14 in levende koolstof toen dezelfde was als in de standaard; dat was niet zo. De kalenderleeftijd wordt afgelezen aan een \omterm{def:b1:radiometric-methods:radiocarbon}{kalibratiecurve} zoals IntCal20, als een reeks jaren cal BP.
\end{solution}

\begin{solution}{exo:b1:radiometric-methods:4}
(a)~Elke klok start wanneer zijn \omterm{def:b1:rock-forming-minerals:mineral}{mineraal} onder zijn \omterm{def:b1:radiometric-methods:closure}{sluitingstemperatuur} afkoelt:
zirkoon (boven \qty{750}{\celsius}) sluit eerst, hoornblende (rond
\qty{530}{\celsius}) daarna, biotiet (rond \qty{280}{\celsius}) als laatste, zodat de leeftijden
in die volgorde afnemen terwijl de graniet afkoelde. (b)~$(530 - 280)/(410 - 398) =
250/12 = \qty{21}{\celsius}$ per miljoen jaar. (c)~Zirkoon sloot boven
\qty{750}{\celsius}, dus koelde de graniet met ten minste $750 - 530 =
\qty{220}{\celsius}$ af in \qty{10}{\Myr}: ten minste \qty{22}{\celsius} per miljoen
jaar, een ondergrens.
\end{solution}

\begin{solution}{exo:b1:radiometric-methods:5}
(a)~$\bar x = 7.5$, $\bar y = 0.809525$, $S_{xx} = 81.0$, $S_{xy} = 1.13965$;
$m = 1.13965/81.0 = 0.0140698$, $b = 0.809525 - 7.5 \times 0.0140698 = 0.70400$.
(b)~$t = \ln(1.0140698)/\num{1.3972e-11} = \qty{1000.0}{\Ma}$; \omterm{def:b1:radiometric-methods:isochron}{beginverhouding}
$0.7040$. (c)~De residuen zijn $+1.6$, $-2.5$, $+0.7$ en $+0.2$ (in eenheden van
\num{e-4}): $s = \sqrt{\num{9.34e-8}/2} = 0.00022$, $\sigma_m = 0.00022/\sqrt{81}
= \num{2.40e-5}$ en $\sigma_t = \num{2.40e-5}/(\num{1.3972e-11} \times 1.0141) =
\qty{1.7}{\Myr}$: \qty[uncertainty-mode=separate]{1000.0(17)}{\Ma}. (d)~$0.7040$
is dicht bij de mantelwaarde $0.7025$ en ver onder de korstwaarde $0.710$: een bron
vooral in de \omterm{def:b1:layered-earth:layers}{mantel}.
\end{solution}

\begin{solution}{exo:b1:radiometric-methods:6}
(a)~$0.100/39.098 \times \num{1.17e-4} = \qty{2.992e-7}{mol}$ kalium-40 per
gram. (b)~${}^{40}\mathrm{Ar}^*/{}^{40}\mathrm{K} = \num{1.169e-9}/\num{2.992e-7}
= 0.0039065$; $\lambda/\lambda_e = 5.543/0.581 = 9.5404$; $t = \ln(1 + 9.5404
\times 0.0039065)/\num{5.543e-10} = \ln(1.037269)/\num{5.543e-10} =
\qty{66.01}{\Ma}$. (c)~Het radiogene argon zou $0.995 \times 1.169$
nanomol zijn, wat \qty{65.69}{\Ma} geeft: de gemeten leeftijd zou \qty{0.32}{\Myr}
te oud zijn, tien keer de \num{32000} jaar waarbinnen de inslag en de
extinctie aan het einde van het Krijt nu gelijktijdig blijken. Overtollig argon is vooral een zorg voor jonge gesteenten en hele gesteenten; een
\omterm{def:b1:rock-forming-minerals:mineral}{mineraal} zoals sanidien zou bij het kristalliseren geen argon moeten opnemen, en een
argon--argonspectrum zou het aan het licht brengen.
\end{solution}

\begin{solution}{exo:b1:radiometric-methods:7}
(a)~Ongeveer \num{2000} cal BP: de curve gaat daar door \num{2070} jaar BP.
(b)~Van ongeveer \num{2495} tot \num{2699} cal BP, tweehonderd jaar: de curve is
bijna vlak tussen \num{2500} en \num{2700} cal BP, zodat elke kalenderleeftijd in dat
interval bijna dezelfde \omterm{def:b1:radiometric-methods:radiocarbon}{radiokoolstofleeftijd} heeft, en een nauwkeuriger meting het
bereik niet zou versmallen. (c)~Bij \num{2000} cal BP is de \omterm{def:b1:radiometric-methods:radiocarbon}{radiokoolstofleeftijd}
\num{70} jaar \emph{ouder} dan de kalenderleeftijd; bij \num{8000} cal BP
\num{822} jaar jonger (\qty{10}{\percent}); bij \num{12000} cal BP \num{1694} jaar jonger (\qty{14}{\percent}). Het verschil is niet evenredig met de leeftijd: het verandert zelfs van teken. Het kan niet alleen van de halveringstijd komen, die een constante \qty{2.4}{\percent} zou geven.
\end{solution}

\begin{solution}{exo:b1:radiometric-methods:8}
(a)~Het asbed eronder (\qty{94.50}{\Ma}) geeft een maximumleeftijd, de lavastroom erboven
(\qty{93.80}{\Ma}) een minimum: de grens is ertussen omsloten. De dijk is
ook een minimum, maar een zwakker, jonger dan de lava. De detritische zirkoon is een
maximum alleen voor het zandsteen, en zwakker dan het asbed, dat onder het
zandsteen ligt. (b)~$f = (12 - 4)/(20 - 4) = 0.5$; $t = 94.50 - 0.5 \times 0.70 =
\qty{94.15}{\Ma}$; $\sigma_t = \sqrt{0.25 \times 0.10^2 + 0.25 \times 0.12^2} =
\qty{0.08}{\Myr}$: \qty[uncertainty-mode=separate]{94.15(8)}{\Ma}. (c)~Haar
korrels kwamen uit een gesteente van ten minste \qty{95.3}{\Ma} oud, geërodeerd ergens in het
brongebied; het toont dat het zandsteen jonger is dan \qty{95.3}{\Ma}, wat het asbed eronder al strakker vastlegt.
\end{solution}

\begin{solution}{exo:b1:radiometric-methods:9}
(a)~Korrel~A: $\ln(1.08401)/\num{1.55125e-10} = \qty{520.0}{\Ma}$ en
$\ln(1.6688)/\num{9.8485e-10} = \qty{520.0}{\Ma}$. Korrel~B: \qty{368.4}{\Ma} en
\qty{389.9}{\Ma}. (b)~Korrel~A is concordant. (c)~$0.05881/0.08401 = 0.7000$ en
$0.4682/0.6688 = 0.7001$. De leeftijd is $\ln(1 + R)/\lambda$, niet evenredig met
$R$: voor de kleine uranium-238-verhouding is $\ln(1+R) \approx R$ en daalt de leeftijd met
bijna \qty{30}{\percent} (tot \qty{368}{\Ma}); voor de grotere uranium-235-verhouding
daalt de logaritme minder, en daalt de leeftijd alleen tot \qty{390}{\Ma}. Beide zijn te
jong, en zij spreken elkaar tegen: de korrel is discordant.
\end{solution}

\begin{solution}{exo:b1:radiometric-methods:10}
(a)~Met $u = \lambda t$, $D^*/P = e^u - 1 = u + u^2/2 + \dots$, dus $(D^*/P)/\lambda
= t\,(1 + u/2 + \dots)$: de benadering overschat de leeftijd met een relatieve fout dicht bij $u/2$. (b)~Rubidium-87 bij \qty{310}{\Ma}: $u = 0.00433$, fout
\qty{0.22}{\percent} (\qty{0.7}{\Myr}); kalium-40 bij \qty{296}{\Ma}: $u =
0.164$, fout \qty{8.7}{\percent} (\qty{26}{\Myr}; $u/2$ geeft \qty{8.2}{\percent});
uranium-238 bij \qty{4.4}{\Ga}: $u = 0.683$, fout \qty{43}{\percent} --- de
benadering en zelfs haar foutschatting ($u/2 = \qty{34}{\percent}$) falen.
(c)~$R = 0.004341$, $\sigma_R = \num{4.34e-6}$; volgens
\cref{prop:b1:radiometric-methods:precision}, $\sigma_t =
\num{4.34e-6}/(\num{1.3972e-11} \times 1.0043) = \qty{0.31}{\Myr}$.
\end{solution}

\begin{solution}{exo:b1:radiometric-methods:11}
(a)~Laat ${}^{40}\mathrm{K}$ het kalium-40 van vandaag zijn. Op het tijdstip $t_2$ vóór het heden waren er ${}^{40}\mathrm{K}\,e^{\lambda t_2}$ atomen, en het argon dat tussen $t_1$ en $t_2$ werd gemaakt was, volgens \cref{prop:b1:radiometric-methods:kar} toegepast op dat interval,
$(\lambda_e/\lambda)\,{}^{40}\mathrm{K}\,e^{\lambda
t_2}\bigl(e^{\lambda (t_1 - t_2)} - 1\bigr) = (\lambda_e/\lambda)\,
{}^{40}\mathrm{K}\,\bigl(e^{\lambda t_1} - e^{\lambda t_2}\bigr)$. Een fractie
$1 - f$ ervan bleef over. Sinds $t_2$ heeft het gesloten \omterm{def:b1:rock-forming-minerals:mineral}{mineraal}
$(\lambda_e/\lambda)\,{}^{40}\mathrm{K}\,(e^{\lambda t_2} - 1)$ meer gemaakt. Optellen en
vermenigvuldigen met $\lambda/(\lambda_e\,{}^{40}\mathrm{K})$ geeft de uitdrukking.
(b)~De schijnbare leeftijd is $\ln(1 + \text{uitdrukking})/\lambda$: \qty{400.0}{\Myr} voor
$f = 0$, \qty{342.3}{\Myr} voor $0.3$, \qty{282.7}{\Myr} voor $0.6$ en
\qty{200.0}{\Myr} voor $1$. (c)~Elke leeftijd tussen \qty{200}{\Myr} en
\qty{400}{\Myr} kan door een gedeeltelijk verlies worden voortgebracht, en één getal kan niet
zeggen of het een werkelijke gebeurtenis op dat tijdstip registreert of een mengsel; een leeftijdsspectrum,
of andere klokken in hetzelfde gesteente, zijn nodig.
\end{solution}

\begin{solution}{exo:b1:radiometric-methods:12}
(a)~De leeftijd is evenredig met de gebruikte gemiddelde levensduur: $10306 \times 8223/8033 =
\num{10550}$ jaar. (b)~$10550 - 10306 = 244$ jaar, ongeveer een zevende van de
\num{1694}. (c)~Als de aanvankelijke activiteit $k A_0$ was, is de ware leeftijd $t = \tau
\ln(k A_0/A) = \num{10550} + \tau \ln k$; met $t = \num{12000}$, $\ln k =
1450/8223 = 0.176$ en $k = 1.19$: ongeveer \qty{19}{\percent} meer koolstof-14 dan
in de standaard.
\end{solution}

\begin{solution}{pb:b1:radiometric-methods:1}
\begin{enumerate}
  \item De punten stijgen gestaag met $x$; hun afwijkingen van een lijn zijn ten hoogste twee eenheden in de vierde decimaal, tegen een spreiding van $0.04$ in $y$:
        zij liggen op een rechte lijn met een kleine spreiding.
  \item $\bar x = 4.90333$, $\bar y = 0.729288$, $S_{xx} = 62.1253$, $S_{xy} =
        0.269335$; $m = 0.269335/62.1253 = 0.0043354$.
  \item $b = 0.729288 - 0.0043354 \times 4.90333 = 0.70803$: de
        $^{87}\mathrm{Sr}/^{86}\mathrm{Sr}$ van de graniet toen zij kristalliseerde, die alle delen van de graniet gemeen hadden.
  \item $t = \ln(1.0043354)/\num{1.3972e-11} = \qty{309.6}{\Ma}$.
  \item De residuen zijn $+9.1$, $-17.2$, $+6.9$, $-6.4$, $+15.1$ en $-7.5$ (in
        eenheden van \num{e-5}); $s = 0.000137$, $\sigma_m = 0.000137/\sqrt{62.1253} =
        \num{1.74e-5}$, $\sigma_t = \num{1.74e-5}/(\num{1.3972e-11} \times 1.0043)
        = \qty{1.2}{\Myr}$: \qty[uncertainty-mode=separate]{309.6(12)}{\Ma}, en
        $\sigma_b = 0.00010$.
  \item $0.7080$ ligt dichter bij de continentale \omterm{def:b1:layered-earth:layers}{korst} $0.710$ dan bij de \omterm{def:b1:layered-earth:layers}{mantel} $0.7025$.
        Voor een mengsel: $(0.70803 - 0.7025)/(0.710 - 0.7025) = 0.74$: ongeveer drie
        kwart van het strontium kwam uit gesmolten \omterm{def:b1:layered-earth:layers}{korst}.
  \item $0.0740/39.098 \times \num{1.17e-4} = \qty{2.214e-7}{mol}$ per gram,
        maal $\num{6.022e23}$: \num{1.33e17} atomen.
  \item $^{40}\mathrm{Ar}^*/^{40}\mathrm{K} = \num{4.139e-9}/\num{2.2144e-7} =
        0.018691$; $1 + 9.5404 \times 0.018691 = 1.17832$; $t =
        \ln(1.17832)/\num{5.543e-10} = \qty{296.0}{\Ma}$.
  \item $^{40}\mathrm{K} = \qty{3.292e-8}{mol.g^{-1}}$, verhouding $0.019239$,
        $\ln(1.18355)/\num{5.543e-10} = \qty{304.0}{\Ma}$.
  \item $\ln(1.018691)/\num{5.543e-10} = \qty{33.4}{\Ma}$, negen keer te jong,
        omdat slechts ongeveer één verval op de tien argon geeft. Calcium-40 is het gewone
        isotoop van calcium, overvloedig in gesteenten: het weinige dat door verval werd gemaakt, kan niet
        worden onderscheiden van het calcium dat de mineralen van het begin af bevatten.
  \item Overtollig argon verschuift leeftijden met bedragen van de orde van een miljoen jaar, wat voor jonge gesteenten telt;
        hier zou het ruim onder \qty{1}{\percent} van de leeftijd blijven. En \omterm{def:b1:rock-forming-minerals:groups}{micas} en \omterm{def:b1:rock-forming-minerals:groups}{amfibool}
        kristallen nemen bijna geen argon op wanneer zij vormen.
  \item Ofwel koelde de graniet langzaam af en sloten de \omterm{def:b1:rock-forming-minerals:groups}{micas} en het \omterm{def:b1:rock-forming-minerals:groups}{amfibool} na de
        kristallisatie (\omterm{def:b1:radiometric-methods:closure}{afkoelingsleeftijden}), ofwel dreef een latere verhitting een deel van hun argon uit.
        Een vlak plateau in een argon--argonleeftijdsspectrum ondersteunt een \omterm{def:b1:radiometric-methods:closure}{afkoelingsleeftijd};
        een spectrum zonder plateau waarschuwt dat de mineralen argon verloren. De volgorde van de leeftijden,
        hoornblende vóór biotiet zoals hun \omterm{def:b1:radiometric-methods:closure}{sluitingstemperaturen} voorspellen, wijst ook op afkoeling.
  \item $\ln(1.049361)/\num{1.55125e-10} = \qty{310.6}{\Ma}$ en
        $\ln(1.35784)/\num{9.8485e-10} = \qty{310.6}{\Ma}$: concordant.
  \item Zirkoon \qty{310.6}{\Ma}, sluiting boven \qty{750}{\celsius}:
        kristallisatie; de \omterm{def:b1:radiometric-methods:isochron}{isochroon} van het hele gesteente,
        \qty[uncertainty-mode=separate]{309.6(12)}{\Ma}, dateert ook
        het gesteente, en komt overeen binnen zijn onzekerheid; hoornblende \qty{304.0}{\Ma}: afkoeling
        door ongeveer \qty{530}{\celsius}; biotiet \qty{296.0}{\Ma}: afkoeling
        door ongeveer \qty{280}{\celsius}. De grafiek loopt van boven
        \qty{750}{\celsius} bij \qty{310.6}{\Ma} naar \qty{280}{\celsius} bij
        \qty{296.0}{\Ma}.
  \item $(530 - 280)/(304.02 - 296.03) = 250/7.99 = \qty{31}{\celsius}$ per
        miljoen jaar.
  \item $(750 - 530)/(310.60 - 304.02) = 220/6.58 = \qty{33}{\celsius}$ per
        miljoen jaar ten minste, omdat de zirkoon boven \qty{750}{\celsius} sloot: de afkoeling
        was in het begin ten minste even snel.
  \item Het temperatuurverschil, $250 \pm \sqrt{40^2 + 40^2} = 250 \pm
        57$, geeft \qtyrange{24}{38}{\celsius} per miljoen jaar.
  \item $(550 - 350)/7.99 = \qty{25}{\celsius}$ per miljoen jaar, binnen het bereik van
        vraag~17: de onzekerheid komt van de \omterm{def:b1:radiometric-methods:closure}{sluitingstemperaturen}, niet van de leeftijden.
  \item $(280 - 10)/30 = \qty{9.0}{km}$.
  \item $\ln(1.039108)/\num{1.55125e-10} = \qty{247.30}{\Ma}$, en hetzelfde uit de
        uranium-235-verhouding: concordant.
  \item Uit het graniet, geërodeerd bij de \omterm{def:b1:cross-sections:unconformity}{discordantie}: hun leeftijd is haar kristallisatieleeftijd. Zij tonen alleen dat het zandsteen jonger is dan
        \qty{310}{\Ma}: een maximumleeftijd, veel zwakker dan het asbed.
  \item $(58 - 18)/(247.30 - 245.30) = \qty{20.0}{m}$ per miljoen jaar; de basis
        is $18/20.0 = \qty{0.90}{\Myr}$ ouder dan het onderste asbed:
        \qty{248.20}{\Ma}.
  \item Van \qty{309.6}{\Ma} (\omterm{def:b1:radiometric-methods:isochron}{isochroon}) of \qty{310.6}{\Ma} (zirkoon) tot
        \qty{248.2}{\Ma}: ongeveer \qty{62}{\Myr}.
  \item \qty{9.0}{km} in $296.0 - 248.2 = \qty{47.8}{\Myr}$: \qty{0.19}{km} per
        miljoen jaar, \qty{0.19}{mm} per jaar. Dit neemt aan dat de gedateerde biotiet
        vlak onder de \omterm{def:b1:cross-sections:unconformity}{discordantie} vandaan kwam, zodat al het gesteente erboven in de graniet en de schist
        was verwijderd tegen de tijd dat het zand arriveerde, en dat de \omterm{def:b1:heat-budget:geotherm}{geotherm} constant bleef.
  \item De graniet koelde af met ongeveer \qty{31}{\celsius} per miljoen jaar
        tussen de sluitingen van hoornblende en biotiet.
\end{enumerate}
\end{solution}
